% !TeX root = ../../tesi.tex

\section{Profilo aziendale e linee di attività}

Halue è una società di consulenza tecnologica nata nel 2022, con sede legale a Parma e sede
operativa ad Abano Terme, in provincia di Padova. È nata come \emph{boutique} di consulenza, cioè
una società piccola e specializzata in pochi ambiti, sull’\gloss{e-commerce}{\emph{e-commerce}} e
sui sistemi di \gloss{crm}{\emph{Customer Relationship Management} (CRM)}, con Shopify e Salesforce
come piattaforme di riferimento. Il primo ambito riguarda la vendita di prodotti e servizi tramite
canali digitali. Il secondo riguarda la gestione delle relazioni delle aziende con i propri
clienti, dalle vendite all’assistenza al \emph{marketing}.

Negli anni successivi Halue ha esteso la propria attività all’intelligenza artificiale, a partire
dai progetti stessi. I negozi dei suoi clienti hanno cataloghi ampi, contenuti in più lingue e
percorsi d’acquisto complessi. In questi punti un modello linguistico collegato ai dati del negozio
aiuta sia chi acquista sia chi gestisce il negozio. Nello stesso periodo le piattaforme di
commercio elettronico hanno iniziato a offrire strumenti e protocolli pensati per gli agenti di
intelligenza artificiale, che presento nella \secref{sec:ecosistema}.

Dal 2024 Halue è una \gloss{societa-benefit}{Società Benefit}: una forma giuridica, introdotta in
Italia nel 2015, per le società che oltre al profitto perseguono obiettivi di beneficio comune
dichiarati nello statuto e ne rendono conto ogni anno in una relazione d’impatto. Gli obiettivi di
Halue riguardano l’ambiente, le persone e le comunità. L’azienda misura le proprie pratiche con il
\gloss{b-impact-assessment}{\emph{B Impact Assessment}}, un questionario di valutazione che le
esamina in cinque aree: \emph{governance}, lavoratori, clienti, comunità e ambiente. Tra i
risultati del 2025, Halue riporta\textsuperscript{\ref{rif:halue-benefit}}:

\begin{itemize}
	\item un codice etico e politiche interne sulla segnalazione di illeciti, sugli abusi e sulle
	discriminazioni;
	\item linee guida per rendere sostenibile il lavoro da remoto;
	\item una collaborazione per la tutela della fauna marina nel Delta del Po.
\end{itemize}

Oggi l’azienda conta dieci persone, con i seguenti ruoli:

\begin{itemize}
        \item \gloss{founder-senior-consultant}{\emph{Founder \& Senior Consultant}}: il fondatore
	Valentino Baraldo, che svolge anche attività di consulenza;
        \item \gloss{solution-technical-architect}{\emph{Solution \& Technical Architect}}:
	definisce le soluzioni e l’architettura tecnica;
        \item \gloss{senior-project-manager}{\emph{Senior Project Manager}} e
	\gloss{project-management-assistant}{\emph{Project Management Assistant}}: due ruoli dedicati alla
	gestione dei progetti;
        \item \gloss{senior-software-engineer}{\emph{Senior Software Engineer}},
	\gloss{software-engineer}{\emph{Software Engineer}} e \gloss{junior-developer}{\emph{Junior
	Developer}}: cinque sviluppatori con diversi livelli di esperienza;
        \item \gloss{ux-ui-digital-consultant}{\emph{UX/UI Digital Consultant}}: cura la
	progettazione dell’esperienza utente e dell’interfaccia.
\end{itemize}

In un’azienda di queste dimensioni ogni ruolo copre un’intera fase del progetto: la consulenza
iniziale, l’architettura, la gestione, lo sviluppo e la progettazione dell’interfaccia. Ogni
persona lavora per più clienti nello stesso periodo.

Halue lavora su commessa: costruisce soluzioni su misura per i singoli clienti e non vende prodotti
propri, come applicazioni pubblicate sullo Shopify App Store. Le sue competenze crescono quindi
progetto dopo progetto, guidate dalle richieste dei clienti, che l’azienda segue da vicino in un
rapporto di lavoro continuo, possibile per un’azienda delle sue dimensioni.

La consulenza tecnologica accompagna tutte le linee di attività dell’azienda, dalla valutazione
iniziale alla scelta dei sistemi. Le linee sono, in ordine di peso:

\begin{enumerate}
        \item l’\emph{e-commerce} su Shopify e su Salesforce Commerce Cloud, sia verso il
	consumatore finale (\gloss{b2c}{B2C}) sia verso altre aziende (\gloss{b2b}{B2B}). Nel B2B il
	negozio gestisce anche listini riservati, ordini minimi e condizioni di pagamento per ciascuna
	azienda cliente;
        \item l’integrazione di sistemi, cioè il collegamento del negozio con la logistica, con i
	servizi di terze parti e con i sistemi gestionali (\gloss{erp}{ERP}), che riuniscono in un’unica
	base di dati amministrazione, acquisti, magazzino e vendite e con cui il negozio sincronizza ordini
	e disponibilità dei prodotti. Il collegamento può essere punto a punto, costruito per ogni coppia di
	sistemi, oppure passare per una piattaforma di integrazione erogata come servizio
	(\gloss{ipaas}{iPaaS}), che offre connettori già pronti e strumenti per trasformare i dati;
        \item i servizi di intelligenza artificiale, dalla valutazione di quanto un’azienda sia
	pronta ad adottarla, anche per la qualità e l’organizzazione dei suoi dati, fino allo sviluppo di
	agenti su misura;
	\item i sistemi CRM su Salesforce, compresi i prodotti per l’assistenza ai clienti e per il
	\emph{marketing};
	\item la formazione, rivolta sia ai clienti sia ai dipendenti. Per i clienti Halue prepara
	incontri come un \emph{workshop} del fondatore sul funzionamento dei modelli linguistici, che
	aggiorna il cliente su questa tecnologia e lo rende consapevole delle sue possibilità e dei suoi
	limiti. I dipendenti seguono una formazione interna sui metodi agili, sulle piattaforme e sulla
	tecnologia, e percorsi personali come il corso di \emph{project management} che un collega
	seguiva durante il mio \emph{stage}: si formano per crescere come professionisti, anche oltre il
	loro ruolo in azienda;
	\item la progettazione dell’interfaccia e dell’esperienza utente.
\end{enumerate}

Le due piattaforme di \emph{e-commerce} hanno pesato in modo diverso nella storia di Halue. Nei
primi anni l’azienda si è specializzata anche su Salesforce Commerce Cloud, la piattaforma di
commercio elettronico della \emph{suite} di Salesforce, e ne segue ancora i clienti. Oggi la
specializzazione principale è Shopify Plus, che presento nella \secref{sec:clientela}. Halue cura
anche la migrazione completa di un negozio da una piattaforma all’altra.

L’intelligenza artificiale è insieme una linea a sé e, sempre di più, una competenza che entra
nei progetti di \emph{e-commerce} e di integrazione.
